\begin{claim}{C.7d}{Proved}
Let $K/\Q_p$ be finite, $F=K\Q_p(\mu_\infty)$, and
$H=G_F$. The following assertions use full geometric fundamental groups.

Every hyperbolic genus-zero curve over a finite extension of $F$ has
injective ordinary point-section map. Moreover, an ordinary point
section cannot equal a tangential section, and tangential sections
at distinct cusps cannot equal each other. The last assertion
distinguishes the underlying cusps, not all choices of tangent vector.

For $U=\mathbf P^1_K\setminus\{0,1,\infty\}$, let
$\Pi=\pi_1(U_{\bar K})$, and let $N_0$ be the intersection of the kernels
of its monodromy permutation representations on all connected finite
etale geometric covers of $U$ whose proper compactifications have genus
zero. Then, for every ordinary $x\in U(K)$,
\[
                 \Pi^{s_x(H)}\ \subseteq\ N_0.                 \tag{C.7d.1}
\]
In particular its image is trivial in every genus-zero Galois-cover
quotient. For any such quotient, the image of the full fixed group
of a tangential section is exactly the image of the corresponding
cusp inertia. There is an explicit such quotient $A_5$, with inertia
orders $2,3,5$.

The subgroup $N_0$ is nontrivial. More precisely, there are infinitely
many nonalternating finite simple groups $S$ such that every epimorphism
$\Pi\to S$ maps $N_0$ onto $S$. Thus (C.7d.1) supplies genuine
nonsolvable-quotient restrictions, while the family of genus-zero
monodromy quotients does not detect the full profinite group.
An explicit choice is $S=\operatorname{PSL}_2(\mathbf F_{47})$,
using the source-qualified known classification stated below.
No nontriviality of the arithmetic fixed group is asserted.
\end{claim}
\begin{proof}
\emph{Fields, paths, and the local Kummer input.}
Every finite extension of $F$ inside $\bar K$ is $M\Q_p(\mu_\infty)$
for some finite extension $M/K$: choose finitely many field generators
over $F$ and adjoin them to $K$. A curve, finitely many points, finite
covers, and tangential parameters over this field all descend to a
finite local extension after enlarging $M$. We may therefore work
over such an $M$ throughout and then restrict to
$H_M=G_{M\Q_p(\mu_\infty)}$.

Recall the proved injection (C.7b.1)
\[
 H^1(G_M,\Q_p(1))\ \hookrightarrow\ H^1(H_M,\Q_p(1)).
                                                               \tag{C.7d.2}
\]
For completeness, the cyclotomic quotient
$\Gamma_M=\operatorname{Gal}(M\Q_p(\mu_\infty)/M)$ is abelian
and its $p$-adic cyclotomic character is nontrivial. If $c$ is a
continuous $1$-cocycle with values in $\Q_p(1)$, commutativity gives
\[
  (\chi(\sigma)-1)c(\tau)=(\chi(\tau)-1)c(\sigma).
\]
Choose $\sigma$ with $\chi(\sigma)\ne1$. The identity makes $c$ a
coboundary, so inflation--restriction proves (C.7d.2).
For $u\in M^\times$, vanishing of its restricted rational Kummer class
therefore says that $u$ is a root of unity. Indeed the valuation part
vanishes, and the local logarithm is injective on an open principal-unit
subgroup, with finite torsion as its kernel on the full unit group.

Equality of two point-section classes is equivalent to the existence
of a Galois-invariant geometric etale path between their fibre functors.
This follows directly from the fundamental-groupoid interpretation:
a geometric conjugator making the two splittings equal is precisely
such a path. For a finite etale cover, a path has a unique lift after
its starting fibre point is specified. If the cover and starting point
are rational over the field in question and the original path is
Galois invariant, uniqueness makes the lifted path invariant as well.
We use this observation after finite local extensions, so every chosen
endpoint can be made rational.

\emph{Ordinary points on genus-zero curves.}
First suppose that $P,Q$ are ordinary points of the tripod over
$M\Q_p(\mu_\infty)$ and that a geometric path between them is
$H_M$-invariant. Enlarge $M$ so that $P,Q\in U(M)$.
Use the finite etale cover
\[
 W=\mathbf P^1\setminus
          \bigl(\{0,\infty\}\cup\mu_3\bigr)\longrightarrow U,
                  \qquad x=z^3.
\]
Choose a starting lift $\alpha$ with $\alpha^3=P$.
The lifted path ends at some $\beta$ with $\beta^3=Q$.
After a further finite extension containing $\alpha,\beta,\mu_3$,
this path is invariant. Each unit $z-\eta$, $\eta\in\mu_3$, pushes
it to an invariant path on $\mathbf G_m$. Consequently
\[
       \kappa\!\left(\frac{\alpha-\eta}{\beta-\eta}\right)
                         \bigm|_{H_M}=0
                 \qquad(\eta\in\mu_3).                       \tag{C.7d.3}
\]
The arguments are nonzero because $P,Q\ne1$.
By (C.7d.2), all three displayed ratios are roots of unity.

The finitely generated characteristic-zero field
$\Q(\alpha,\beta,\mu_3)$ embeds into $\C$; no assertion that its
generators are algebraic over $\Q$ is needed. In this embedding,
(C.7d.3) says that $\alpha$ and $\beta$ have equal distances to the
three cube roots of unity. These centres are noncollinear. Subtracting
two squared-distance identities gives two independent real linear
equations for $\alpha-\beta$, so $\alpha=\beta$. Hence $P=Q$.

For a general hyperbolic genus-zero curve, extend the finite local
field so that its compactification is $\mathbf P^1$ and three of its
cusps are rational. Fill in all other cusps and identify the resulting
curve with the tripod. Distinct ordinary points remain distinct.
Their equal sections would contradict the tripod case. This proves
the first point-section assertion over every finite extension of $F$.

\emph{Ordinary points and tangential cusps.}
Suppose an invariant path joins a tangent at $0$ to an ordinary
point $t$ of the tripod. Lift it through $z^3=x$. The starting
tangent lifts to a tangent at $z=0$, and the other endpoint is an
ordinary $\alpha$ with $\alpha^3=t$. After adjoining the necessary
cube root of the leading tangent coefficient, the lift is rational
and invariant over a finite local extension. Explicitly, a Puiseux
parameter with $x=a u$ lifts by $z=a^{1/3}u^{1/3}$; the cyclotomic
Galois group fixes all roots of unity, and the chosen leading
coefficient has been made rational.

Now fill in $z=0$. On the hyperbolic genus-zero curve
\[
          W^\circ=\mathbf P^1\setminus(\mu_3\cup\{\infty\}),
\]
the starting section becomes the ordinary point $0$ and the other
one remains the ordinary point $\alpha\ne0$. The preceding
injectivity is a contradiction. By changing coordinates and
retaining three cusps, this distinguishes an ordinary point from a
tangential section on every hyperbolic genus-zero curve.

If an invariant path instead joins tangents at two distinct cusps,
first retain three cusps including them, with the first at $0$.
Lift through $z^3=x$ and fill in $z=0$ as above. The first section
becomes ordinary and the other is still tangential at a point of
$\mu_3$ or at infinity. The assertion just proved gives a contradiction.
This argument makes no claim that distinct leading tangent
coefficients at the same cusp always give distinct restricted sections.

\emph{Finite monodromy representations.}
Let $Y\to U_{\bar K}$ be a connected finite etale cover whose proper
compactification has genus zero, and let $g\in\Pi^{s_x(H)}$.
Choose a finite local extension $M/K$ over which the cover and all
points of the fibre over $x$ are defined. Start the lift of the loop
$g$ at an arbitrary such point $P$. Its endpoint $Q$ is again in
that fibre. The lifted path is $H_M$-invariant, so ordinary
point-section injectivity on $Y$ gives $P=Q$. The curve $Y$ is
hyperbolic: its negative Euler characteristic is the degree of the
cover times that of $U$. Since every starting point is fixed,
$g$ acts trivially on the entire fibre. This proves (C.7d.1), including
non-Galois covers and their possibly nonsolvable monodromy groups.

For a Galois cover with group $G$, apply the same lifting argument
to a full fixed loop based at a tangent. After a finite local
extension its starting and ending fibre functors are tangential at
cusps of $Y$. The cusp distinction just proved makes their underlying
cusps equal. To identify the resulting group constraint precisely,
label the geometric fibre by $G$, with deck transformations acting on
the left and geometric monodromy on the right. For the selected
base cusp, let $I\le G$ be the image of its inertia. The cusps
above it correspond to the right $I$-orbits in $G$. A loop with
monodromy image $h$ lifts from the fibre label $1$ to $h$; equality
of their underlying cusps says $h\in I$. Thus the fixed-group
image is contained in $I$. The full cusp inertia in $\Pi$ is
itself fixed by the tangential section over $H$, because its
conjugation action is the cyclotomic character, which is trivial
on $H$. It maps onto $I$, proving equality.

\emph{An explicit nonsolvable quotient.}
Let $\zeta$ be a primitive fifth root of unity, and put
\[
 a=\zeta-\zeta^4,\qquad b=\zeta^2-\zeta^3,\qquad
 S=\begin{pmatrix}a&b\\b&-a\end{pmatrix},\qquad
 T=\begin{pmatrix}\zeta^2&0\\0&\zeta^3\end{pmatrix}.
                                                               \tag{C.7d.4}
\]
In $\operatorname{PGL}_2(\Q(\zeta))$ these generate $A_5$.
Here is a finite exact verification, also recorded in the independent
standard-library certificate
\texttt{code/BC\_genus\_zero\_detection\_checks.py}.
Work in $\Q[z]/(z^4+z^3+z^2+z+1)$ and normalize a projective
matrix by dividing by its first nonzero entry. One checks
$a^2+b^2=-5$, $(ab)^2=5$, and projective orders
$|S|=2$, $|T|=5$, $|ST|=3$. Closing the two normalized
matrices under right multiplication produces exactly sixty matrices.
Their element orders have multiplicities
\[
  \begin{array}{c|rrrr}
    \text{order}&1&2&3&5\\ \hline
    \text{multiplicity}&1&15&20&24.
  \end{array}
\]
The centralizers of the fifteen involutions are exactly five distinct
Klein-four groups. Conjugation on those five groups is faithful
and its image consists of all sixty even permutations. For the
certificate's labelling of the five objects, the generators act by
\[
     S=(2\ 4)(3\ 5),\qquad T=(1\ 5\ 3\ 4\ 2).
                                                               \tag{C.7d.5}
\]
The certificate verifies the conjugation action for every product,
not merely the generator orders. This supplies the claimed group
identification entirely by rational arithmetic.

The quotient of $\mathbf P^1$ by this finite group has genus zero
by Luroth's theorem. In characteristic zero a finite point stabilizer
in $\operatorname{PGL}_2$ is cyclic: its action on the tangent line
is injective, since a nontrivial finite-order automorphism with
derivative one cannot fix a point. The order list above therefore
shows that the fifteen order-two, ten order-three, and six order-five
cyclic subgroups are the complete stabilizers at their respective
fixed points. Each subgroup fixes two points, and different such
stabilizers cannot share a point. There are consequently $30,20,12$
ramified points of indices $2,3,5$. These are single group orbits,
of lengths $60/2,60/3,60/5$, so the quotient has exactly three
branch values. Normalize these to $0,1,\infty$ after a finite
extension of the field of definition. Removing their inverse images
gives the required genus-zero $A_5$ Galois cover of the tripod.
The Riemann--Hurwitz check is
$-120+30(2-1)+20(3-1)+12(5-1)=-2$.
All of this takes place in characteristic zero, so the construction
works over finite extensions of every $\Q_p$, including $p=2,3,5$.
In particular one can place the order-five branch value at $0$:
the ordinary full fixed image in this $A_5$ quotient is $1$,
while the tangential full fixed image is its designated $C_5$ inertia.

\emph{A proved limit on the detection family.}
Guralnick's Theorem~1.1, with the definitions immediately preceding it,
states that for each fixed genus there are only finitely many
nonalternating nonabelian simple composition factors of monodromy
groups of tame covers with source of at most that genus
\cite[pp.~2--3, Theorem~1.1]{BC_Guralnick2003}.
In characteristic zero every finite cover here is tame. Let
$\mathcal S_0$ be the finite set supplied by that theorem for genus zero.
Its scope concerns composition factors of all monodromy groups, so
no assumption that our rational covers are primitive is being made.

There are infinitely many two-generated nonalternating finite simple
groups outside $\mathcal S_0$: for example take
$\operatorname{PSL}_2(\mathbf F_q)$ for sufficiently large prime $q$.
The upper and lower elementary unipotent matrices generate
$\operatorname{SL}_2(\mathbf F_q)$, hence their images generate its
simple projective quotient. Since $\Pi$ is free profinite of rank two,
it surjects onto any such $S$. Fix an epimorphism $\varphi:\Pi\to S$.

There is also an explicit choice, using an additional \emph{Known}
classification input. Burness's opened author proceedings version
defines $E(g)$ in composition-factor terms on author page~7 and
in Section~4.3, author page~36, gives the complete list
\[
 \operatorname{PSL}_2(\mathbf F_q)\in E(0)
 \quad\Longleftrightarrow\quad
 q\in\{7,8,11,13,16,17,19,25,27,29,31,32,37,41,43,64\},
\]
where the alternating isomorphism cases have been excluded by definition
\cite[author pp.~7,36, Section~4.3]{BC_Burness2018}.
In particular $S=\operatorname{PSL}_2(\mathbf F_{47})$ is excluded.
The source attributes this list to Frohardt--Guralnick--Magaard (2002).
We opened the stated list and its composition-factor definition, but
could not open that original classification proof; this is a cited
known input, not a reproof of the classification. A concrete epimorphism
from $\Pi$ sends its two free generators to the projective classes of
\[
       \begin{pmatrix}1&1\\0&1\end{pmatrix},\qquad
       \begin{pmatrix}1&0\\1&1\end{pmatrix}
                           \quad\hbox{over }\mathbf F_{47}.
\]
Their powers give all upper and lower elementary matrices, so Gaussian
elimination proves that they generate; the certificate additionally
enumerates all $47(47^2-1)/2=51888$ projective classes.

Suppose that $N_0\subseteq\ker\varphi$. The complement of
$\ker\varphi$ is compact. By the definition of $N_0$, the complements
of the genus-zero monodromy kernels cover that compact set.
Finitely many suffice. Thus, for finitely many corresponding
surjections $\rho_i:\Pi\to G_i$,
\[
                  \bigcap_i\ker\rho_i\ \subseteq\ker\varphi.
\]
The image $L$ of $\Pi$ in $\prod_iG_i$ is a subdirect product and
surjects onto $S$. Every composition factor of a finite subdirect
product is a composition factor of one of its factors. To see this
without assuming the same assertion for arbitrary subgroups, project
$L$ onto its first $r-1$ factors. Its kernel is a subgroup of the last
factor $G_r$ that is normal in $G_r$, because the projection
$L\to G_r$ is surjective. Induction and a composition series prove
the assertion.

Consequently $S$ would be a composition factor of some $G_i$,
contradicting its choice. We have proved
$N_0\not\subseteq\ker\varphi$. Since $N_0$ is normal in $\Pi$
and $S$ is simple, $\varphi(N_0)=S$. This proves the asserted
nontriviality and the stronger quotient statement. It identifies
the exact limit of (C.7d.1) as a detection method. C.7h below
determines the arithmetic fixed subgroup by a different argument;
that subgroup must not be replaced by the much larger subgroup $N_0$.
\end{proof}
